% SentinelEdge -- EMNLP 2026 System Demonstrations
% Demo session description (1-page supplement)
%
% This document is intended as a companion to the main paper. It is
% NOT part of the main paper's page count. EMNLP demos publish only
% the main paper; this supplement is for the demo session organizers
% and reviewers who want to evaluate what the live demonstration will
% look like.

\documentclass[11pt,a4paper]{article}
\usepackage[T1]{fontenc}
\usepackage[utf8]{inputenc}
\usepackage[margin=2.2cm]{geometry}
\usepackage{microtype}
\usepackage{times}
\usepackage{graphicx}
\usepackage{enumitem}
\usepackage{url}
\usepackage{titling}
\usepackage[hidelinks]{hyperref}
\setlength{\droptitle}{-3em}

\title{\textbf{SentinelEdge: Live Demonstration Description}\\
{\large EMNLP 2026 System Demonstrations -- supplement}}
\author{}
\date{}

\begin{document}
\maketitle
\vspace{-3em}

\paragraph{What visitors will see.}
The demo is a single browser window (Figure~\ref{fig:layout}) connected
to a local backend by WebSocket. Visitors are offered a panel of nine
pre-loaded call transcripts (six scam categories: IRS impersonation,
tech support, bank fraud, Amazon refund, crypto investment, prize
notification, plus the grandparent and utility-shutoff scams; three
legitimate calls: doctor appointment, pharmacy refill, airline
rebooking). They may also paste an arbitrary transcript and replay
it. When a call is started, the backend streams the transcript one
sentence at a time with realistic $1.5$--$3$ second pauses, scoring
each sentence through the trained model in the live code path. The UI
renders, in real time:

\begin{itemize}[leftmargin=*,itemsep=2pt,topsep=2pt]
\item \textbf{Phone-call UI} -- a phone simulator showing the
incoming call, the caller's claimed identity, and the streaming
transcript as it accumulates.
\item \textbf{Risk gauge} -- an animated dial showing the current EMA
score against the four threshold zones (safe, low, medium, high,
critical). When the EMA crosses $0.75$, a red alert overlay appears
with an explanation.
\item \textbf{Feature breakdown} -- a list of which of the $18$
handcrafted features fired on the most recent sentence. For example,
on the sentence ``\textit{This is an arrest warrant from the IRS}'' the
panel highlights \texttt{impersonation\_count = 1},
\texttt{has\_threat = 1}, and \texttt{urgency\_count = 0}.
\item \textbf{Alert explanation} -- human-readable reasons generated
by the alert engine, e.g. \textit{``urgency language detected;
impersonation pattern: \texttt{irs}''}.
\item \textbf{Latency display} -- the measured wall-clock latency of
the most recent sentence's classification, taken from
\texttt{time.perf\_counter()} around the actual
\texttt{predict\_proba} call.
\end{itemize}

\paragraph{What runs locally.}
The entire system runs on the demo laptop. The backend is a single
Python process (\texttt{demo/backend/main.py}) running on
\texttt{localhost:8000}. The frontend is a Vite-built React app
served from \texttt{localhost:5173}. No network calls leave the
machine during a demonstration. A console panel below the main UI
exposes the in-process logs so visitors can verify this claim by
inspection.

\paragraph{Hardware and setup.}
The demo runs comfortably on any modern laptop with at least 4 GB of
free RAM and Python 3.10+. We will bring our own laptop (MacBook Pro,
M-series) and a backup. Setup takes under five minutes: install
dependencies, run the training pipeline once
(\texttt{python training/generate\_synthetic\_data.py \&\& python
training/prepare\_datasets.py \&\& python training/fit\_tfidf.py \&\&
python training/train\_call\_classifier.py}), then start the backend
and frontend. The repo ships an idempotent script
(\texttt{deploy\_server.py}) that performs the full setup in one
command. A poster summarizing the architecture, the headline numbers
from the main paper, and a printed QR code linking to the
reproducibility scripts will accompany the demo.

\paragraph{Interactive elements.}
Visitors can: (i) pick any of the pre-loaded calls and watch the EMA
trajectory build sentence-by-sentence; (ii) interrupt a call mid-stream
via a ``block'' button to see the alert pipeline tear down; (iii)
paste arbitrary transcripts (e.g., examples brought from their own
research) and watch the model classify them in real time; (iv) toggle
between the default and adversarially-retrained classifier checkpoint
to see the differing scores on \textit{polite-social-engineering}
examples. A short walk-through demonstration (under three minutes per
visitor) covers all four interactions.

\paragraph{What the demo deliberately does not claim.}
The Whisper transcription stage, the live microphone capture, the
Android \texttt{CallScreeningService} integration, and the federated
learning loop are described in the main paper but \emph{not} exercised
during the live demonstration. The demo's input is text only. This
matches Section~7 of the main paper and is signposted on the demo
poster so that visitors know exactly which parts of the pipeline they
are seeing live.

\paragraph{Failure modes and contingencies.}
The demo has been tested in three failure modes: (i) backend crash
(frontend displays a reconnect prompt and the visitor can restart);
(ii) missing model artefacts (the patched backend now refuses to start
with a clear error message rather than silently falling back to a
heuristic scorer); (iii) WebSocket disconnect (the frontend buffers UI
state for $10$ seconds before clearing). The full pipeline can also
be exercised offline from the command line via the reproducibility
scripts, so if the live demo cannot start, we can fall back to a
recorded walkthrough.

\begin{figure}[ht]
\centering
\fbox{\parbox{0.92\textwidth}{\centering\footnotesize\vspace{4pt}
\textbf{[ DEMO SCREEN LAYOUT -- TO BE REPLACED WITH SCREENSHOT ]}\\[6pt]
{\renewcommand{\arraystretch}{1.0}
\begin{tabular}{c@{\quad}c}
\textit{Phone simulator + live transcript} & \textit{Risk gauge + alert overlay} \\[2pt]
\textit{Feature breakdown panel} & \textit{Alert reasons + latency display} \\
\end{tabular}}\vspace{4pt}
}}
\caption{Browser layout of the live demonstration. Left: streaming
transcript inside a phone-call UI. Right: real-time risk gauge,
feature breakdown, alert reasons, and per-sentence inference
latency. Replace this placeholder with a real screenshot at
camera-ready.}
\label{fig:layout}
\end{figure}

\paragraph{Reproducibility from the demo.}
A QR code on the demo poster links to a public anonymized repository
containing the trained models, evaluation harness, and a one-command
reproducibility script (\texttt{python experiments/run\_all.py}). A
visitor who copies the repository can reproduce every number in the
main paper, including the time-to-detection CDF, the latency Pareto,
and the ASR-robustness curves, on a single CPU core in roughly ten
minutes.

\end{document}
